\part{Application thorn writing}

%%%%%%%%%%%%%%%%%%%%%%%%%%%%%%%%%%%%%%%%%%%%%%%%%%%%%%%%%%%%%%%%%%%%%%%
%%%%%%%%%%%%%%%%%%%%%%%%%%%%%%%%%%%%%%%%%%%%%%%%%%%%%%%%%%%%%%%%%%%%%%%

\chapter{Thorn concepts} 

      (This section has to contain enough explanation to make the rest of
       the writers guide readable the first time through)
(The hardest bugs to find are
those arising from plausible but incorrect assumptions about the
behavior of someone else's thorn.)
      a) Again probably emphasize collaboration, what are thorns,
         toolkits, how to share them.
      b) Things to think about before you start programming:
             Language, read all the documentation, emphasize use of
             standard supported Cactus infrastructure
      c) Available data types
         i)   Scalars
         ii)  Arrays and GFs
         iii) Groups
      d) Ghost zones and parallelism
      e) Understanding the RFR concept
      f) Understanding the GH concept

%%%%%%%%%%%%%%%%%%%%%%%%%%%%%%%%%%%%%%%%%%%%%%%%%%%%%%%%%%%%%%%%%%%%%%%
%%%%%%%%%%%%%%%%%%%%%%%%%%%%%%%%%%%%%%%%%%%%%%%%%%%%%%%%%%%%%%%%%%%%%%%

\chapter{Anatomy of a thorn}

Each thorn must be in a toolkit. Toolkits are fairly informal. If you
want, you can just make your own toolkit directory, and stick your new
thorn into it.

      a) Creating a thorn (gmake new thorn)
      b) Thorn directory structure
      c) Thorn make files
      d) Thorn configuration
         (i)   Parameters (param.ccl)
         (ii)  Variables (interface.ccl)
         (iii) Scheduler (schedule.ccl)

%%%%%%%%%%%%%%%%%%%%%%%%%%%%%%%%%%%%%%%%%%%%%%%%%%%%%%%%%%%%%%%%%%%%%%%
%%%%%%%%%%%%%%%%%%%%%%%%%%%%%%%%%%%%%%%%%%%%%%%%%%%%%%%%%%%%%%%%%%%%%%%

\chapter{Putting code into your thorn}

      a) What the flesh basically provides (header files, macros, ...)
      b) Argument lists and parameters
      c) A simple example (Baloney)
      d) Differences between F and C with the interface
      e) A more complex example (Wavetoy)
      f) Error handling and code stopping
      g) Calling C from F, F from C, (C++ from F? ....)
      h) Good Cactus Style Programming Guidelines
      i) Adding a test suite
      j) Sharing your thorn with the world.

%%%%%%%%%%%%%%%%%%%%%%%%%%%%%%%%%%%%%%%%%%%%%%%%%%%%%%%%%%%%%%%%%%%%%%%

\appendix{Appendix}

\chapter{Using GNATS}
\chapter{Using CVS}
\chapter{Using TAGS}
    \section{TAGS with emacs}
    \section{TAGS with vi}


\chapter{Configuration file syntax}
    
\section{General Concepts}

Each thorn is configured by three compulsory files in the top level thorn
directory:
\begin{itemize}
\item{} interface.ccl
\item{} param.ccl
\item{} schedule.ccl
\end{itemize}

\section{interface.ccl}

The interface configuration file consists of a header block giving details 
of the thorns relationship with other thorns, and then a series of blocks 
listing the thorn's global variables.
The header block has the form:
{\t 
\begin{verbatim}
implements: <implementation>
[inherits: <implementation>, <implementation>]
[friend: <implentation>, <implementation>]
\end{verbatim}
}
where
\begin{itemize}
\item{} The implentation name must be unique among all thorns, unless ????.
\item{} Inheriting from another implentation makes all that implentations public 
        variables and parameters available to your thorn. At least one thorn 
        providing any inherited implentation must be present at compile time.
        A thorn cannot inherit from itself. Inheritance is associative, but not
        commutative.
\item{} Being a friend of another implementation makes all that implementations 
        protected and public (????)
        variables and parameters available to your thorn. A friend is not required 
        for compilation. A thorn cannot be its own friend. Friendship is 
        associative and commutative.
\end{itemize} 

The thorn's variables are defined by:
{\t
\begin{verbatim}
[<access>:]

VTYPE <group name> [TYPE=<gtype>] [DIM=<dim>] 
[{
   <variable name>[,]<variable name>
   <variable name>
}]
\end{verbatim}}

\begin{itemize}
\item{} {\t access} defines which thorns have can use the following 
        groups, and can be either {\t public}, {\t protected} or {\t private}.
\item{} {\t VTYPE} defines the data type of the variables in the group.
Supported data types are {\t LOGICAL}, {\t INTEGER}, {\t CHAR} and {\t REAL}.
\item{} {\t group name} must be a alpha-numeric name which is unique 
within the scope of the thorn. A group name is compulsory.
\item{} {\t TYPE} designates the kind of variables held by the group.
The choices are {\t GF}, {\t ARRAY} or {\t SCALAR}. This field is 
optional, with the default variable type being {\t SCALAR}.
\item{} {\t DIM} defines the spatial dimension if the group is of type
{\t ARRAY} or {\t GF}, and can take the value {\t 1}, {\t 2}, or {\t 3}.
The default value is {\t DIM=3}.
\item{} The block following the group declaration line contains a list of
        variables contained in the group. All variables in a group have 
        the same data type, variable type and dimension. The list can be 
        separated by spaces, commas, or new lines. The variable names 
        must be unique within the scope of the thorn. 
        A variable can only be a member of 
        one group. The block must be delimited by brackets on new lines.
        If no block is given after a group declaration line, a 
        variable with the same name as the group is created.
\end{itemize} 

\section{param.ccl}

\section{schedule.ccl}

The schedule configuration files consists of 
\begin{itemize}
\item{} assignments to switch on storage and communications 
        for array variables at the start of program execution.
\item{} {\it schedule blocks} which schedule a subroutine from the
        thorn to be called at a specified time during program 
        executation. Statements within the schedule block can
        be used to switch on storage and communication for groups
        of variables during the duration that the subroutine is called.
\item{} Conditional statements
\end{itemize}

{\it assignments statements} have the form:
{\t
\begin{verbatim}
[STORAGE: <group>, <group>]
[COMMUNICATION: <group>, <group>]
\end{verbatim}
}

Each {\it schedule block} in the file {\t schedule.ccl} must have
the syntax:
{\t
\begin{verbatim}
schedule [<function name>] at <time>
{
  [LANG: <language>]
  [STORAGE: <group>,<group>]
  [COMMUNICATION: <group>,<group>]
  [TRIGGER: <grid function>,<grid function>]
}
\end{verbatim}
}

Any schedule block or assignment statements can be optionally 
surrounded by conditional {\t if-elseif-else}
constructs using the parameter data base. These can be nested,
and have the general form:
{\t
\begin{verbatim}
if (CCTK_Equals(<parameter>,<string>)) 
{
  [<assignments>]
  [<schedule blocks>]
}
\end{verbatim}
Conditional constructs cannot be used inside of a schedule block.


\chapter{CCTK Parameters}
\chapter{Thorn-Flesh Interface}

%%%%%%%%%%%%%%%%%%%%%%%%%%%%%%%%%%%%%%%%%%%%%%%%%%%%%%%%%%%%%%%%%%%%%%%
%%%%%%%%%%%%%%%%%%%%%%%%%%%%%%%%%%%%%%%%%%%%%%%%%%%%%%%%%%%%%%%%%%%%%%%
%%%%%%%%%%%%%%%%%%%%%%%%%%%%%%%%%%%%%%%%%%%%%%%%%%%%%%%%%%%%%%%%%%%%%%%

\end{document}
